\documentclass[11pt]{article}
\usepackage{notes}
\definecolor{AccN}{HTML}{6B2D8B}
\begin{document}
\titlepg{3}{Supervised Learning Part 2:\\Regression and Sensor Fusion}{OLS, Ridge, Lasso, Elastic Net, SVR, gradient boosting, prediction uncertainty, the Kalman filter and ML-enhanced fusion}
\tableofcontents
\newpage

\section{Regression and the Two Case Studies}
\begin{definition}{regression}
Learning a function $f:\mathbb R^d\to\mathbb R$ from examples $(x_i,y_i)$ with real-valued targets, e.g.\ mission duration in hours, fuel consumption, casualties, or target position.
\end{definition}
\textbf{Case A} predicts time-to-mission-completion from pre-mission parameters: $N=6{,}200$ exercise records (70/15/15 split), 18 features including terrain, weather, force size, armour count, air assets, OPFOR density, EW intensity, supply level, mission type, rules of engagement, distance and elevation change. Targets range from 2.1 to 48.6 hours. \textbf{Case B} tracks five airborne targets by fusing three sensors: ground radar (range, bearing, speed; 1 Hz), infrared (bearing, elevation; 2 Hz; blocked by cloud 20\% of the time) and ELINT (frequency, PRF, bearing; 5 Hz).

\begin{pitfall}{outliers in regression}
Squared-error methods are pulled strongly by extreme cases --- one 48-hour siege among two-hour raids. Plot residuals, report RMSE together with $R^2$, and consider robust losses.
\end{pitfall}

\section{Ordinary Least Squares}
\subsection{Model and closed-form solution}
The linear model is $y=Xw+\varepsilon$ with design matrix $X\in\mathbb R^{N\times d}$ and errors $\varepsilon\sim\mathcal N(0,\sigma^2I)$.
\begin{derivation}{the normal equations}
Minimise $L(w)=\|y-Xw\|^2=(y-Xw)^\top(y-Xw)$. Expanding, $L=y^\top y-2w^\top X^\top y+w^\top X^\top Xw$. Setting the gradient to zero,
\[\nabla_wL=-2X^\top y+2X^\top Xw=0\ \Rightarrow\ X^\top Xw=X^\top y\ \Rightarrow\ \boxed{w^*=(X^\top X)^{-1}X^\top y}.\]
The Hessian $2X^\top X$ is positive semi-definite, so this is a minimum; it is unique when $X^\top X$ is invertible (no exact collinearity).
\end{derivation}
\textbf{Geometric view.} The fitted values $\hat y=Xw^*$ are the orthogonal projection of $y$ onto the column space of $X$; the residual $y-\hat y$ is perpendicular to every column, which is exactly the statement $X^\top(y-Xw^*)=0$.

\begin{center}\begin{tikzpicture}
\fill[AccN!12,draw=AccN](0,0)--(5,0)--(6.5,1.4)--(1.5,1.4)--cycle;\node[AccN,font=\small]at(3.2,0.35){column space of $X$};
\draw[-{Stealth},thick,NavyN](1.9,0.55)--(4.3,2.75)node[above right]{$y$};\draw[-{Stealth},thick,GrnN](1.9,0.55)--(4.3,0.62)node[right]{$\hat y=Xw^*$};\draw[dashed,RedN,thick](4.3,0.62)--(4.3,2.75)node[midway,right,font=\small]{residual $\perp$};
\end{tikzpicture}\\*{\small\color{GreyN}Figure 1: least squares as orthogonal projection.}\end{center}

\textbf{Gauss--Markov theorem.} If $E[y|X]=Xw$, errors have equal variance and are uncorrelated, then OLS has the smallest variance among all linear unbiased estimators (it is BLUE). Normal errors are needed only for exact confidence intervals.

\subsection{Case A results and interpretation}
OLS: training $R^2=0.738$, test $R^2=0.710$, RMSE 4.22 h. Coefficients have direct meaning: each EW-intensity level adds 2.31 h, each unit of OPFOR density 1.84 h, each 10 km of distance 0.94 h; each supply grade saves 1.12 h and each 100 troops 0.73 h.

\subsection{Multicollinearity}
When features are nearly linear combinations of each other, $X^\top X$ is nearly singular and coefficients become unstable. The \textbf{variance inflation factor} $\text{VIF}_j=1/(1-R_j^2)$, where $R_j^2$ is from regressing feature $j$ on the others, measures this; force size and armour count had VIF 8.4 (values above about 5--10 are a warning).

\section{Regularised Regression}
\subsection{Ridge (L2)}
\[L_{\text{ridge}}=\|y-Xw\|^2+\lambda\|w\|^2\quad\Rightarrow\quad w^*_{\text{ridge}}=(X^\top X+\lambda I)^{-1}X^\top y.\]
Adding $\lambda I$ makes the matrix invertible for any $\lambda>0$ and shrinks coefficients toward zero smoothly; none become exactly zero. Bias increases slightly while variance falls. Case A ($\lambda=0.8$ by 5-fold CV): test $R^2=0.738$, RMSE 3.91 h, and the force--armour VIF fell to 2.1.

\subsection{Lasso (L1) and why it produces zeros}
\[L_{\text{lasso}}=\|y-Xw\|^2+\lambda\|w\|_1.\]
\begin{derivation}{soft-thresholding in one dimension}
For a single standardised feature with OLS coefficient $z$, minimise $(z-w)^2+\lambda|w|$. For $w>0$ the derivative $-2(z-w)+\lambda=0$ gives $w=z-\lambda/2$, valid only if $z>\lambda/2$; similarly for $w<0$. Otherwise the minimum is at $w=0$. Hence
\[w_{\text{lasso}}=\text{sign}(z)\max\big(|z|-\lambda/2,\;0\big):\]
coefficients smaller than $\lambda/2$ in magnitude are set exactly to zero. Ridge instead gives $w=z/(1+\lambda)$, which shrinks but never reaches zero.
\end{derivation}
Geometrically, the L1 constraint region is a diamond whose corners lie on the axes; the elliptical loss contours usually touch it at a corner, where some coefficients are zero.
\begin{center}\begin{tikzpicture}[x=0.8cm,y=0.8cm]
\begin{scope}\draw[->](-2.2,0)--(2.4,0)node[right]{$w_1$};\draw[->](0,-2)--(0,2.2)node[above]{$w_2$};\draw[BluN,thick](0,0)circle(1.2);\foreach \r in {0.5,1.0,1.5}{\draw[RedN!60,rotate around={35:(1.7,1.4)}](1.7,1.4)ellipse({\r*1.1} and {\r*0.55});}\fill[RedN](0.95,0.74)circle(0.07);\node[font=\small]at(0,-2.5){ridge: circle};\end{scope}
\begin{scope}[xshift=6cm]\draw[->](-2.2,0)--(2.4,0)node[right]{$w_1$};\draw[->](0,-2)--(0,2.2)node[above]{$w_2$};\draw[BluN,thick](-1.2,0)--(0,1.2)--(1.2,0)--(0,-1.2)--cycle;\foreach \r in {0.5,1.0,1.5}{\draw[RedN!60,rotate around={35:(1.7,1.4)}](1.7,1.4)ellipse({\r*1.1} and {\r*0.55});}\fill[RedN](0,1.2)circle(0.07);\node[font=\small]at(0,-2.5){lasso: diamond, solution at $w_1=0$};\end{scope}
\end{tikzpicture}\\*{\small\color{GreyN}Figure 2: why L1 gives sparse solutions.}\end{center}
Case A ($\lambda=0.3$): $R^2=0.728$, 6 of 18 features zeroed (comms quality, temperature, days since resupply, morale, fog of war, enemy experience).

\subsection{Elastic net}
$L_{\text{EN}}=\|y-Xw\|^2+\lambda_1\|w\|_1+\lambda_2\|w\|^2$ combines sparsity with stability among correlated features (lasso alone tends to pick one of a correlated group arbitrarily). Case A: $R^2=0.741$ with 10 features. Rule of thumb: ridge for many correlated relevant features; lasso when most features are irrelevant; elastic net when unsure.

\section{Support Vector Regression}
SVR ignores errors smaller than $\varepsilon$ and penalises larger ones linearly:
\[L_\varepsilon(y,\hat y)=\max(0,|y-\hat y|-\varepsilon),\qquad\min\tfrac12\|w\|^2+C\sum_i(\xi_i+\xi_i^*)\]
subject to $y_i-w^\top x_i-b\le\varepsilon+\xi_i$ and $w^\top x_i+b-y_i\le\varepsilon+\xi_i^*$. Points inside the ``$\varepsilon$-tube'' cost nothing; only points on or outside it become support vectors, and the prediction is $\hat y=\sum_i(\alpha_i-\alpha_i^*)K(x_i,x)+b$. Linear (not squared) penalties make SVR robust to outliers, and kernels capture non-linear effects such as the multiplicative interaction of mountainous terrain and rain.
\begin{center}\begin{tikzpicture}\begin{axis}[width=10cm,height=4.5cm,xlabel={$x$},ylabel={$y$},domain=0:6,samples=60,ticks=none,ymin=-0.8,ymax=3.4]
\addplot[name path=u,draw=none]{1.6+sin(deg(x))*1.1+0.4};\addplot[name path=l,draw=none]{1.6+sin(deg(x))*1.1-0.4};\addplot[AccN!15]fill between[of=u and l];
\addplot[AccN,thick]{1.6+sin(deg(x))*1.1};\addplot[only marks,mark=*,mark size=1.3pt,BluN]coordinates{(0.4,1.9)(1.0,2.6)(1.8,2.5)(2.4,1.9)(3.0,1.6)(3.6,0.9)(4.2,0.4)(5.0,0.6)(5.6,1.5)};\addplot[only marks,mark=*,mark size=1.6pt,RedN]coordinates{(1.4,3.1)(4.6,-0.4)(2.9,2.4)};
\end{axis}\end{tikzpicture}\\*{\small\color{GreyN}Figure 3: the $\varepsilon$-tube; red points outside it incur slack.}\end{center}
Case A ($C=50$, $\varepsilon=0.5$, $\gamma=0.05$): $R^2=0.812$, RMSE 3.24 h, 1,842 support vectors (42.4\% of training data), training 4.8 s. Training cost grows as $O(N^2)$--$O(N^3)$, so for tens of thousands of samples gradient boosting is preferred.

\section{Gradient Boosting for Regression}
Start from $F_0=\bar y$; at each round fit a regression tree $T_m$ to the residuals $r_{im}=y_i-F_{m-1}(x_i)$ (the negative gradient of squared loss) and update $F_m=F_{m-1}+\eta T_m$. XGBoost minimises $\sum_i(y_i-\hat y_i)^2+\sum_m\Omega(T_m)$ with $\Omega(T)=\gamma|T|+\tfrac12\lambda\|w\|^2$, using a second-order approximation for speed. Validation RMSE fell from 6.81 h (10 trees) to 2.11 h (400 trees) and rose again beyond, indicating overfitting --- the number of trees is itself a hyper-parameter, best chosen by early stopping. Final test: $R^2=0.912$, RMSE 2.06 h, MAE 1.54 h; the most important features (EW intensity 0.231, OPFOR density 0.188, distance 0.152) agree with the OLS coefficients.

\section{Regression Metrics and Their Traps}
\[\text{MSE}=\tfrac1N\sum(y_i-\hat y_i)^2,\ \ \text{RMSE}=\sqrt{\text{MSE}},\ \ \text{MAE}=\tfrac1N\sum|y_i-\hat y_i|,\ \ R^2=1-\frac{SS_{\text{res}}}{SS_{\text{tot}}}.\]
RMSE is in the target's units and penalises large errors; MAE is robust; $R^2$ is the share of variance explained; adjusted $R^2=1-(1-R^2)\frac{N-1}{N-d-1}$ penalises extra features. RMSE noticeably larger than MAE (2.06 vs 1.54) signals a few large errors. Traps: compare training and test $R^2$; RMSE hides directional bias (report the signed mean error, since late missions matter more than early ones); $R^2$ is negative when a model is worse than predicting the mean; autocorrelated time-series data inflate $R^2$ unless tested out-of-time; MAPE is undefined when $y=0$.

\section{Prediction Uncertainty}
A point forecast (``18.4 hours'') is incomplete without its uncertainty. \textbf{Aleatoric} uncertainty is inherent randomness (enemy behaviour, weather); \textbf{epistemic} uncertainty reflects lack of data and shrinks with more relevant data. The OLS prediction interval is
\[\hat y_{\text{new}}\pm t_{N-d-1,\alpha/2}\,\hat s\sqrt{1+x^\top(X^\top X)^{-1}x},\]
which widens far from the training data. Tree ensembles give an approximate interval from the spread of individual tree predictions ($\pm3.8$ h average for the 90\% band in Case A). Bayesian ridge regression places a prior $w\sim\mathcal N(0,\tfrac{\sigma^2}{\lambda}I)$, yielding the posterior $\mathcal N\big(w_{\text{ridge}},\,\sigma^2(X^\top X+\lambda I)^{-1}\big)$; its 90\% intervals covered 89.4\% of test cases (well calibrated) and widened to $\pm12$ h or more for unusual missions --- ``18 hours, but possibly 30'' is far more useful to a planner than false precision.

\section{The Kalman Filter}
\subsection{State-space model}
\[x_k=Fx_{k-1}+Bu_k+w_k,\ \ w_k\sim\mathcal N(0,Q);\qquad z_k=Hx_k+v_k,\ \ v_k\sim\mathcal N(0,R).\]
$x_k$ is the hidden state (e.g.\ position and velocity), $F$ the dynamics, $Q$ the process noise, $z_k$ the measurement, $H$ the measurement model and $R$ the sensor noise covariance.

\subsection{Predict and update}
\textbf{Predict:} $\hat x_k^-=F\hat x_{k-1}^++Bu_k$, \ $P_k^-=FP_{k-1}^+F^\top+Q$.\\
\textbf{Update:} innovation $\nu_k=z_k-H\hat x_k^-$, \ $S_k=HP_k^-H^\top+R$, \ $K_k=P_k^-H^\top S_k^{-1}$,
\[\hat x_k^+=\hat x_k^-+K_k\nu_k,\qquad P_k^+=(I-K_kH)P_k^-.\]
\begin{derivation}{the scalar Kalman gain as optimal weighting}
With a prior estimate $\hat x^-$ of variance $P$ and a measurement $z$ of variance $R$ (and $H=1$), consider the combination $\hat x=\hat x^-+K(z-\hat x^-)$. Its error variance is $(1-K)^2P+K^2R$. Differentiating with respect to $K$ and setting to zero gives
\[K=\frac{P}{P+R},\qquad\text{variance}=(1-K)P=\frac{PR}{P+R}<\min(P,R).\]
The filter weights each source inversely to its uncertainty, and the fused estimate is always more certain than either source alone.
\end{derivation}
\begin{example}{one update}
Prior $\hat x^-=100$ km, $P=4$; radar measures 104 km with $R=1$. $K=4/5=0.8$; $\hat x^+=100+0.8(4)=103.2$ km; $P^+=0.2\times4=0.8$.
\end{example}
\begin{center}\begin{tikzpicture}\begin{axis}[width=9cm,height=4.2cm,xlabel={position (km)},ticks=none,domain=93:110,samples=80,legend pos=north west]
\addplot[BluN,thick]{exp(-(x-100)^2/8)/sqrt(8*pi)};\addlegendentry{prior}\addplot[RedN,thick]{exp(-(x-104)^2/2)/sqrt(2*pi)};\addlegendentry{measurement}\addplot[GrnN,thick]{exp(-(x-103.2)^2/1.6)/sqrt(1.6*pi)};\addlegendentry{fused}
\end{axis}\end{tikzpicture}\\*{\small\color{GreyN}Figure 4: Kalman fusion narrows uncertainty.}\end{center}
When $R\gg P$ the gain tends to zero and the filter trusts its prediction; when $R\ll P$ it trusts the measurement. In Case B the state per target is $[X,Y,\dot X,\dot Y]$; radar noise $R=\text{diag}(50^2\,\text{m}^2,0.5^2\,\text{deg}^2,10^2\,(\text{km/h})^2)$, IR is angle-only, and ELINT bearing has $\pm5^\circ$.

\subsection{Non-linear variants}
The \textbf{extended Kalman filter} linearises non-linear dynamics $f$ and measurement models $h$ with Jacobians at the current estimate; it fails when linearisation error is large. The \textbf{unscented Kalman filter} propagates a small set of carefully chosen sigma points through the exact non-linear function and recomputes mean and covariance --- more accurate for manoeuvring or hypersonic targets, with no Jacobians.

\section{Multi-Sensor Fusion}
\textbf{Early (feature-level) fusion} concatenates all sensors' data into one model: it can learn cross-sensor correlations but breaks when a sensor is missing and must handle different sampling rates (Case B: 91.4\%, but 74.8\% with 20\% IR gaps). \textbf{Late (decision-level) fusion} runs a model per sensor and combines decisions: robust to sensor loss but blind to cross-sensor correlations (89.7\%). \textbf{Hybrid fusion} uses a Kalman filter for physics and timing and a learned model to correct its systematic errors: $\hat x_{\text{final}}=\hat x_{KF}+\delta x$, with $\delta x$ predicted by a random forest from residual features (innovation, its squared norm, innovation covariance, time since each sensor's last update, sensor availability). This scored best: 96.2\%, and still 94.8\% with 20\% IR loss, at 2.3 ms latency. Combat violates the Kalman filter's assumptions of linear dynamics and Gaussian noise (evasive manoeuvres, jamming spikes); the learned correction captures what the physics model misses.

\subsection{Missing data and late measurements}
Sensors drop out: radar in terrain shadow or under jamming (30--80\% loss), IR in cloud (15--40\%), ELINT under emission control. In Case B the accuracy loss from different imputation strategies was 2.1 points for mean imputation, 1.4 for carrying the last observation forward, 0.8 for using the Kalman prediction, and 0.3 for ML imputation from the other sensors. Imputed values must be flagged and given inflated measurement noise $R$, otherwise the filter treats them as genuine and becomes falsely confident. \textbf{Out-of-sequence measurements} --- a radar return arriving 0.3 s after later data --- must be applied at their true time by rolling back the filter and re-propagating.

\section{Deeper Dive I: Regression Diagnostics}
\subsection{Reading residual plots}
After fitting, plot residuals $e_i=y_i-\hat y_i$ against fitted values and against each feature. A healthy model shows a structureless horizontal band. Patterns reveal violations:
\begin{itemize}
\item \textbf{Curvature} --- a missing non-linear term (e.g.\ duration rising faster than linearly with distance).
\item \textbf{Funnel shape} (heteroscedasticity) --- error variance grows with the prediction; long missions are less predictable. Remedies: log-transform the target, weighted least squares, or quantile models.
\item \textbf{Isolated large residuals} --- outliers or data errors.
\end{itemize}
\begin{center}\begin{tikzpicture}
\begin{axis}[name=a,width=5.4cm,height=4cm,title={healthy},xlabel={fitted},ylabel={residual},ticks=none,ymin=-3,ymax=3]
\addplot[only marks,mark=*,mark size=1pt,BluN]coordinates{(1,0.4)(2,-0.6)(3,0.8)(4,-0.3)(5,0.5)(6,-0.9)(7,0.2)(8,0.7)(9,-0.5)(10,0.1)(2.5,0.9)(5.5,-0.4)(7.5,-0.7)(9.5,0.6)};\addplot[GreyN,dashed,domain=0:11]{0};
\end{axis}
\begin{axis}[at={(a.east)},anchor=west,xshift=0.8cm,width=5.4cm,height=4cm,title={funnel},xlabel={fitted},ticks=none,ymin=-3,ymax=3]
\addplot[only marks,mark=*,mark size=1pt,RedN]coordinates{(1,0.1)(2,-0.2)(3,0.4)(4,-0.5)(5,0.8)(6,-1.1)(7,1.4)(8,-1.7)(9,2.2)(10,-2.5)(2.5,0.2)(5.5,-0.9)(7.5,1.6)(9.5,-2.2)};\addplot[GreyN,dashed,domain=0:11]{0};
\end{axis}\end{tikzpicture}\\*{\small\color{GreyN}Figure 5: residual plots --- a structureless band versus heteroscedasticity.}\end{center}

\subsection{Leverage and influence}
A point far from the others in feature space has high \textbf{leverage} $h_{ii}=x_i^\top(X^\top X)^{-1}x_i$; if it also has a large residual it is \textbf{influential}, measured by Cook's distance. A single 48-hour siege recorded with an unusual combination of features can swing several coefficients. Always inspect high-leverage cases before trusting coefficients --- the regression analogue of Anscombe's fourth dataset.

\subsection{Interaction and polynomial terms}
Linear regression is linear in the \emph{parameters}, not necessarily in the features. Adding $x_j^2$ or $x_jx_k$ terms captures curvature and interactions: $\text{duration}=\beta_0+\beta_1\,\text{mountain}+\beta_2\,\text{rain}+\beta_3(\text{mountain}\times\text{rain})$, where $\beta_3>0$ means rain hurts more in mountains. With 18 features there are 153 pairwise interactions --- a strong case for regularisation or tree-based models that find interactions automatically.

\section{Deeper Dive II: Choosing $\lambda$ and the Bias--Variance Path}
Regularisation strength is chosen by cross-validation over a logarithmic grid (e.g.\ $\lambda=10^{-3},\ldots,10^{3}$). Validation error typically falls as $\lambda$ removes variance, then rises as bias dominates. The \textbf{one-standard-error rule} picks the largest $\lambda$ whose CV error is within one standard error of the minimum, giving a simpler model that is statistically indistinguishable from the best.
\begin{example}{ridge in one dimension}
With a single standardised feature, OLS coefficient $z=2.4$ and $\lambda=0.8$: ridge gives $2.4/1.8=1.33$; lasso gives $2.4-0.4=2.0$. For a weak feature with $z=0.3$: ridge $0.17$, lasso $0$. Ridge shrinks proportionally; lasso shrinks by a fixed amount and deletes weak effects.
\end{example}
Features must be standardised before penalised regression; otherwise the penalty depends on arbitrary units (a coefficient per metre is penalised differently from one per kilometre).

\section{Deeper Dive III: Quantile Regression and Planning Margins}
Least squares predicts the conditional \emph{mean}. Logistics planners often need an upper quantile: ``a duration we will not exceed in 90\% of missions''. \textbf{Quantile regression} minimises the pinball loss
\[L_\tau(e)=\begin{cases}\tau e & e\ge0\\(\tau-1)e & e<0\end{cases},\qquad e=y-\hat y,\]
which for $\tau=0.9$ penalises under-prediction nine times as much as over-prediction. Gradient-boosting libraries support this loss directly. Planning on the 90th percentile builds the margin into the forecast instead of adding an arbitrary buffer.

\section{Deeper Dive IV: The Kalman Filter in Two Dimensions}
\subsection{A constant-velocity tracker}
State $x=[p,\,v]^\top$ (position, velocity), time step $\Delta t=1$ s:
\[F=\begin{pmatrix}1&\Delta t\\0&1\end{pmatrix},\quad H=\begin{pmatrix}1&0\end{pmatrix},\quad Q=q\begin{pmatrix}\Delta t^3/3&\Delta t^2/2\\\Delta t^2/2&\Delta t\end{pmatrix}.\]
The radar measures position only, yet the filter estimates velocity because successive position innovations reveal it through the off-diagonal terms of $P$.
\begin{example}{one predict--update cycle}
Prior $\hat x=(1000\text{ m},\,200\text{ m/s})$, $P=\text{diag}(100,25)$, $q$ small (take $Q=0$ for clarity), radar $R=400$ m$^2$. \textbf{Predict:} $\hat x^-=(1200,200)$; $P^-=FPF^\top=\begin{pmatrix}125&25\\25&25\end{pmatrix}$. \textbf{Measure} $z=1230$: innovation $\nu=30$, $S=125+400=525$. \textbf{Gain} $K=P^-H^\top/S=(0.238,\,0.048)$. \textbf{Update:} position $1200+0.238\times30=1207.1$ m; velocity $200+0.048\times30=201.4$ m/s. The position innovation also corrected the velocity estimate.
\end{example}

\subsection{Tuning $Q$ and $R$, and checking consistency}
$R$ comes from sensor specifications; $Q$ expresses how much the target may deviate from the motion model. Too small a $Q$ makes the filter sluggish and lose manoeuvring targets; too large makes it noisy. A well-tuned filter has innovations that are zero-mean, white, and of the size it predicts: the \textbf{normalised innovation squared} $\nu^\top S^{-1}\nu$ should average the measurement dimension (1 here). Persistent values above that indicate over-confidence --- exactly the signal the hybrid model's residual corrector learns from.

\subsection{Data association}
With five targets, each radar return must be assigned to a track. \textbf{Gating} accepts a return for a track only if $\nu^\top S^{-1}\nu$ is below a $\chi^2$ threshold; nearest-neighbour, global-nearest-neighbour (Hungarian algorithm) or probabilistic data association resolve conflicts. Mis-association is a major failure mode in dense raids and when an adversary uses decoys.

\section{Deeper Dive V: Information Fusion and Covariance Intersection}
For independent Gaussian estimates the \textbf{information form} is simplest: information $I=P^{-1}$ adds, $I_{\text{fused}}=\sum_iP_i^{-1}$, and $\hat x_{\text{fused}}=I_{\text{fused}}^{-1}\sum_iP_i^{-1}\hat x_i$. If the sensors' errors are correlated (e.g.\ two tracks sharing an earlier common measurement), naive fusion double-counts information and becomes over-confident. \textbf{Covariance intersection} fuses with $P^{-1}=\omega P_1^{-1}+(1-\omega)P_2^{-1}$, choosing $\omega$ to minimise the fused uncertainty --- conservative but safe when correlations are unknown, as in networked air pictures.

\section{Deeper Dive VI: Matrix Regression Worked End to End}
\begin{example}{OLS and ridge with two centred features}
Four missions, centred features (EW intensity, OPFOR density) and centred durations:
\[X=\begin{pmatrix}-1&-1\\-1&1\\1&-1\\1&1\end{pmatrix},\qquad y=\begin{pmatrix}-5\\-1\\1\\5\end{pmatrix}.\]
Then $X^\top X=\begin{pmatrix}4&0\\0&4\end{pmatrix}$ (orthogonal features) and $X^\top y=(-1\cdot-5-1\cdot-1+1\cdot1+1\cdot5,\ \ -1\cdot-5+1\cdot-1-1\cdot1+1\cdot5)=(12,\,8)$.\\
\textbf{OLS:} $w^*=(X^\top X)^{-1}X^\top y=(3,\,2)$: each EW level adds 3 h, each OPFOR unit 2 h (in centred units). Fitted values $(-5,-1,1,5)$ reproduce $y$ exactly, $R^2=1$.\\
\textbf{Ridge,} $\lambda=4$: $(X^\top X+4I)^{-1}X^\top y=\tfrac18(12,8)=(1.5,\,1.0)$ --- both coefficients halved. With orthogonal features ridge shrinks every coefficient by the factor $d/(d+\lambda)$ where $d$ is the feature's sum of squares.
\end{example}

\section{Deeper Dive VII: A Bayesian Prediction Interval}
\begin{example}{how uncertain is a forecast?}
A ridge-type Bayesian model predicts 18.4 h for a new mission, with noise variance $\sigma^2=9$ h$^2$ (s.d.\ 3 h) and posterior variance of the mean prediction $x^\top\Sigma_wx=4$ h$^2$ for a typical mission. The predictive variance is $9+4=13$, s.d.\ $3.6$ h; a 90\% interval is $18.4\pm1.645\times3.6=(12.5,\,24.3)$ h. For an unusual mountain--jungle mission far from the training data, $x^\top\Sigma_wx$ might be 45: predictive s.d.\ $\sqrt{54}=7.3$ h and the 90\% interval widens to $(6.3,\,30.5)$ h. Same point forecast, very different planning implications --- the epistemic term grows where data are sparse.
\end{example}

\clearpage
\part*{Part C --- Further Theory, Methods and Applications}
\addcontentsline{toc}{part}{Part C --- Further Theory, Methods and Applications}

\section{Regression as a Statistical Model}
\subsection{Least squares is maximum likelihood}
Assume $y_i=w^\top x_i+\varepsilon_i$ with independent Gaussian noise $\varepsilon_i\sim\mathcal N(0,\sigma^2)$. The log-likelihood of the data is
\[\ln L(w,\sigma)=-\frac N2\ln(2\pi\sigma^2)-\frac1{2\sigma^2}\sum_i(y_i-w^\top x_i)^2,\]
so maximising the likelihood over $w$ is exactly minimising the sum of squared residuals. Squared error is therefore not arbitrary: it is the right loss when errors are Gaussian. When they are heavy-tailed (occasional wild outliers), the absolute or Huber loss corresponds to a better noise model.

\subsection{The Gauss--Markov theorem}
If the errors have zero mean, equal variance and no correlation (Gaussianity is not needed), OLS has the smallest variance among all \emph{linear unbiased} estimators: it is BLUE. Two consequences: when the assumptions hold, no unbiased linear rival does better; and the way to beat OLS is to accept some bias --- which is exactly what ridge and lasso do.

\subsection{How certain are the coefficients?}
Under the Gaussian model each coefficient has a standard error, and the ratio of estimate to standard error follows a $t$ distribution with $N-p$ degrees of freedom ($p$ parameters). For a single predictor,
\[\hat\sigma^2=\frac{SS_{\text{res}}}{N-2},\qquad SE(\hat b_1)=\frac{\hat\sigma}{\sqrt{S_{xx}}},\qquad S_{xx}=\sum_i(x_i-\bar x)^2.\]
\begin{example}{inference for the EW-intensity model}
Five missions: EW level $x=1,\ldots,5$; duration $y=6,9,9,13,13$ h. OLS gives $\hat y=4.6+1.8x$ with $SS_{\text{res}}=3.6$. Then $\hat\sigma^2=3.6/3=1.2$, $\hat\sigma=1.095$; $S_{xx}=10$; $SE(\hat b_1)=1.095/\sqrt{10}=0.346$; $t=1.8/0.346=5.20$ with 3 degrees of freedom, two-sided $p=0.014$. The 95\% confidence interval for the slope is $1.8\pm3.182\times0.346=(0.70,\;2.90)$ hours per EW level. The effect is clearly positive, but with five missions its size is uncertain by a factor of four.
\end{example}

\subsection{Confidence intervals versus prediction intervals}
Two different questions have two different intervals at a new input $x_0$:
\[\text{mean response: }\hat y_0\pm t\,\hat\sigma\sqrt{\tfrac1N+\tfrac{(x_0-\bar x)^2}{S_{xx}}},\qquad\text{one new mission: }\hat y_0\pm t\,\hat\sigma\sqrt{1+\tfrac1N+\tfrac{(x_0-\bar x)^2}{S_{xx}}}.\]
\begin{example}{how long will the next mission take?}
At EW level 4: $\hat y_0=11.8$ h. The 95\% interval for the \emph{average} duration of level-4 missions is $11.8\pm1.9$ h; the 95\% \emph{prediction} interval for one particular mission is $11.8\pm4.0$ h. At level 8 --- outside the data --- the prediction is 19.0 h with a prediction interval of $\pm6.7$ h, and even that assumes the straight line continues. A planner needs the prediction interval, and needs to know when $x_0$ is an extrapolation.
\end{example}
\begin{center}\begin{tikzpicture}\begin{axis}[width=8cm,height=5cm,xmin=0,xmax=9,ymin=0,ymax=30,xlabel={EW level $x$},ylabel={duration (h)},grid=both,grid style={GreyN!15},font=\scriptsize,legend style={font=\tiny,at={(0.02,0.98)},anchor=north west}]
\addplot[forget plot,name path=pu,draw=none,domain=0:9,samples=50]{4.6+1.8*x+3.182*1.0954*sqrt(1+0.2+(x-3)^2/10)};
\addplot[forget plot,name path=pl,draw=none,domain=0:9,samples=50]{4.6+1.8*x-3.182*1.0954*sqrt(1+0.2+(x-3)^2/10)};
\addplot[AccN!15,area legend]fill between[of=pu and pl];\addlegendentry{95\% prediction band}
\addplot[forget plot,name path=cu,draw=none,domain=0:9,samples=50]{4.6+1.8*x+3.182*1.0954*sqrt(0.2+(x-3)^2/10)};
\addplot[forget plot,name path=cl,draw=none,domain=0:9,samples=50]{4.6+1.8*x-3.182*1.0954*sqrt(0.2+(x-3)^2/10)};
\addplot[BluN!25,area legend]fill between[of=cu and cl];\addlegendentry{95\% confidence band}
\addplot[NavyN,thick,domain=0:9]{4.6+1.8*x};\addlegendentry{fitted line}
\addplot[only marks,mark=*,mark size=1.6pt]coordinates{(1,6)(2,9)(3,9)(4,13)(5,13)};
\draw[dashed,GreyN](axis cs:5.5,0)--(axis cs:5.5,30);\node[font=\tiny,GreyN,anchor=west]at(axis cs:5.6,3){extrapolation};
\end{axis}\end{tikzpicture}\\*{\small\color{GreyN}Figure C1: both bands widen away from the centre of the data; the prediction band is always wider.}\end{center}

\section{Counting Casualties: Generalised Linear Models}
Casualty numbers, equipment failures and intrusion attempts are \textbf{counts}: non-negative integers whose spread grows with their mean. Ordinary regression can predict negative counts and assumes constant variance; the natural model is \textbf{Poisson regression}, a member of the family of generalised linear models (GLMs).
\subsection{Poisson regression}
\[Y\sim\text{Poisson}(\mu),\qquad\ln\mu=\beta_0+\beta_1x_1+\cdots+\beta_px_p.\]
The log link keeps $\mu$ positive and makes effects multiplicative: increasing $x_j$ by one unit multiplies the expected count by the \textbf{rate ratio} $e^{\beta_j}$. When units are observed for different lengths of time or with different strengths, an \textbf{offset} $\ln(\text{exposure})$ is added to the linear predictor, so the model predicts rates per unit of exposure.
\begin{example}{a casualty-estimate model}
From simulated engagements, a planning cell fits expected friendly casualties per company-day as $\ln\mu=0.5+0.25\,x$, where $x$ is the number of enemy artillery battalions in range. Rate ratio $e^{0.25}=1.28$: each extra battalion raises expected casualties by 28\%. Predictions: $x=0$: $\mu=1.65$; $x=2$: $\mu=2.72$; $x=4$: $\mu=4.48$. Because counts are Poisson, the planner can also state probabilities: at $x=4$, $P(\text{no casualties})=e^{-4.48}=1.1\%$ and $P(\ge8\text{ casualties})=8.5\%$. Medical and evacuation capacity should be planned against the upper tail, not the mean.
\end{example}
\subsection{Over-dispersion}
In the Poisson model the variance equals the mean. Real casualty data are usually \textbf{over-dispersed}: a few engagements are catastrophic. Check by comparing the residual deviance (or Pearson $\chi^2$) with its degrees of freedom; a ratio well above 1 indicates over-dispersion. The \textbf{negative binomial} model adds a dispersion parameter, $\text{Var}(Y)=\mu+\mu^2/k$, and gives wider, more honest intervals. Ignoring over-dispersion produces standard errors that are too small and confidence intervals that are falsely precise.
\subsection{Other GLMs}
Binary outcomes (mission success, target destroyed) use logistic regression --- the binomial GLM; positive skewed quantities (repair times, costs) use gamma regression with a log link; proportions of a known total use binomial regression. The recipe is always the same: choose a distribution that matches the data, a link function that keeps predictions in range, and fit by maximum likelihood.
\begin{pitfall}{casualty estimates are planning inputs, not forecasts of fact}
Casualty-estimation models are fitted to simulations and historical data that differ from the coming operation in force mix, terrain, medical support and enemy behaviour. Report them as ranges with stated assumptions, validate them against independent data where any exist, and remember that they inform decisions about human lives: the uncertainty must be visible to the commander, not hidden in a point estimate.
\end{pitfall}

\section{Force Effectiveness: Regression Meets Lanchester}
\subsection{Lanchester's equations}
Lanchester's \textbf{square law} for aimed fire models the strengths $B(t)$ and $R(t)$ of two forces by
\[\frac{dB}{dt}=-r\,R,\qquad\frac{dR}{dt}=-b\,B,\]
where $r$ and $b$ are the per-unit kill rates (effectiveness coefficients). Dividing and integrating gives the invariant $b(B_0^2-B^2)=r(R_0^2-R^2)$: fighting strength is proportional to effectiveness times the \emph{square} of numbers, so concentration pays. The \textbf{linear law}, for unaimed area fire, has $dB/dt=-r\,R\,B$, and strength scales linearly with numbers.
\subsection{Estimating effectiveness coefficients by regression}
The coefficients are rarely known; they are estimated from simulation runs or historical engagements. Under the square law, Blue's loss rate is proportional to Red's strength, so $r$ is the slope of a regression through the origin of Blue losses per hour on Red strength:
\[\hat r=\frac{\sum_tR_t\,\Delta B_t}{\sum_tR_t^2}.\]
\begin{example}{estimating Red's kill rate}
Three hourly observations from a simulated engagement: Red strength 100, 80, 60 tanks; Blue losses 5.1, 4.0, 3.1 tanks per hour. $\sum R\,\Delta B=510+320+186=1{,}016$; $\sum R^2=20{,}000$; $\hat r=0.0508$ Blue tanks killed per Red tank per hour. The fitted rates $5.08,4.06,3.05$ match the observations closely, which supports the square-law form for this engagement.
\end{example}
\subsection{Testing which law applies}
A more general model is $\Delta B_t=c\,R_t^{\,p}B_t^{\,q}$. Taking logarithms turns it into a linear regression,
\[\ln\Delta B_t=\ln c+p\ln R_t+q\ln B_t+\varepsilon_t,\]
whose fitted exponents are tested against the square law ($p=1,q=0$) and the linear law ($p=1,q=1$). Studies of historical battles have found exponents between the two, varying with terrain and era --- a reminder that the laws are idealisations. Modern force-effectiveness models replace the single coefficient by a regression on terrain, weather, morale, electronic warfare and training: in the Lecture~1 example, a random forest with these factors raised explained variance of engagement outcomes from $R^2=0.61$ to $0.87$ over 1,200 engagements.
\begin{defence}{force-ratio planning}
Once $r$ and $b$ are estimated, the square law gives the force ratio needed for Blue to prevail: Blue wins if $bB_0^2>rR_0^2$, i.e.\ $B_0/R_0>\sqrt{r/b}$. If Red's kill rate is four times Blue's, Blue needs twice Red's numbers, not four times. Because $\sqrt{r/b}$ is a non-linear function of two uncertain estimates, its uncertainty should be propagated (for example by bootstrapping the engagements) before it is used in planning.
\end{defence}

\section{Flexible Regression and Honest Uncertainty}
\subsection{Non-parametric regression}
\begin{itemize}
\item \textbf{$k$-NN regression} predicts the average of the $k$ nearest training targets.
\item \textbf{Kernel (Nadaraya--Watson) regression} averages all targets weighted by a kernel of distance, $\hat y(x)=\sum_iK(x,x_i)y_i/\sum_iK(x,x_i)$.
\item \textbf{Splines} fit piecewise polynomials joined smoothly at knots; smoothing splines penalise curvature.
\item \textbf{Regression trees and forests} (Lecture~2) fit piecewise-constant surfaces.
\end{itemize}
All of them adapt to non-linear relationships without a formula; all of them extrapolate poorly beyond the data.

\subsection{Gaussian process regression}
A Gaussian process (GP) places a probability distribution over functions: any finite set of function values is jointly Gaussian with covariance given by a kernel, typically $k(x,x')=\sigma_f^2\exp\big(-(x-x')^2/2\ell^2\big)$. Conditioning on training data gives, at a new point $x_*$,
\[\mu_*=k_*^\top(K+\sigma_n^2I)^{-1}y,\qquad\sigma_*^2=k(x_*,x_*)-k_*^\top(K+\sigma_n^2I)^{-1}k_*,\]
where $K$ is the kernel matrix of the training inputs and $k_*$ the vector of kernels between $x_*$ and each training input.
\begin{example}{a GP with two observations}
Observations $y=1$ at $x=0$ and $y=2$ at $x=1$; $\ell=1$, $\sigma_f=1$, noise variance 0.01. At $x_*=0.5$ (between the data): $\mu_*=1.64$, standard deviation 0.19 --- confident. At $x_*=3$ (far from the data): $\mu_*=0.29$, standard deviation 0.99 --- the prediction reverts toward the prior mean of zero and the model reports that it does not know. This built-in honesty about extrapolation is why GPs are used for expensive simulations (surrogate models of wargame outcomes) and for choosing which simulation to run next.
\end{example}
\begin{center}\begin{tikzpicture}\begin{axis}[width=8cm,height=4.6cm,xmin=-1.5,xmax=4,ymin=-2.5,ymax=3.5,grid=both,grid style={GreyN!15},font=\scriptsize,xlabel={$x$}]
\addplot[name path=u,draw=none,domain=-1.5:4,samples=80]{(1.01*exp(-x^2/2)-0.60653*exp(-(x-1)^2/2))*1.0/0.6521*1+(1.01*exp(-(x-1)^2/2)-0.60653*exp(-x^2/2))*2/0.6521*1+2*sqrt(max(0.0001,1-((1.01*(exp(-x^2/2))^2-2*0.60653*exp(-x^2/2)*exp(-(x-1)^2/2)+1.01*(exp(-(x-1)^2/2))^2)/0.6521)))};
\addplot[name path=l,draw=none,domain=-1.5:4,samples=80]{(1.01*exp(-x^2/2)-0.60653*exp(-(x-1)^2/2))*1.0/0.6521*1+(1.01*exp(-(x-1)^2/2)-0.60653*exp(-x^2/2))*2/0.6521*1-2*sqrt(max(0.0001,1-((1.01*(exp(-x^2/2))^2-2*0.60653*exp(-x^2/2)*exp(-(x-1)^2/2)+1.01*(exp(-(x-1)^2/2))^2)/0.6521)))};
\addplot[BluN!20]fill between[of=u and l];
\addplot[NavyN,thick,domain=-1.5:4,samples=80]{(1.01*exp(-x^2/2)-0.60653*exp(-(x-1)^2/2))*1.0/0.6521+(1.01*exp(-(x-1)^2/2)-0.60653*exp(-x^2/2))*2/0.6521};
\addplot[only marks,mark=*,AccN]coordinates{(0,1)(1,2)};
\end{axis}\end{tikzpicture}\\*{\small\color{GreyN}Figure C2: GP mean and $\pm2$ standard-deviation band; uncertainty collapses near data and grows away from it.}\end{center}

\subsection{Conformal prediction: intervals without distributional assumptions}
Split conformal prediction wraps \emph{any} regression model in prediction intervals with guaranteed coverage. Fit the model on a training set; compute absolute residuals $|y_i-\hat y_i|$ on a separate calibration set of size $n$; take the $\lceil(n+1)(1-\alpha)\rceil$-th smallest residual as $q$; report $\hat y\pm q$. If future data are exchangeable with the calibration data, the interval contains the truth with probability at least $1-\alpha$.
\begin{example}{an 80\% conformal interval}
Nine calibration residuals (h), sorted: 0.2, 0.4, 0.7, 0.9, 1.1, 1.3, 1.6, 2.3, 3.0. For $\alpha=0.2$: $\lceil10\times0.8\rceil=8$, so $q=2.3$ h. A new mission predicted at 14.0 h receives the interval $(11.7,\;16.3)$ h, whatever model produced the 14.0.
\end{example}
The guarantee fails under distribution shift --- the same caveat as every other method --- but conformal intervals are simple, model-agnostic and increasingly used to attach honest uncertainty to machine-learning predictions.

\section{Forecasting Over Time}
Logistics demand (fuel, ammunition, spares) is a time series: today's value depends on yesterday's. \textbf{Simple exponential smoothing} forecasts the next value as a weighted average that discounts the past geometrically:
\[\hat y_{t+1}=\alpha y_t+(1-\alpha)\hat y_t,\qquad0<\alpha\le1.\]
\begin{example}{forecasting daily fuel demand}
Demand over five days (thousand litres): 120, 135, 128, 150, 142; $\alpha=0.3$, initial forecast 120. Successive one-day-ahead forecasts: 120.0, 124.5, 125.6, 132.9, and for day~6, 135.6. The forecast lags a rising trend; \textbf{Holt's method} adds a smoothed trend term, and \textbf{Holt--Winters} adds seasonality (weekly training cycles).
\end{example}
\textbf{ARIMA} models describe a series by its own lagged values (autoregression), lagged forecast errors (moving average) and differencing to remove trends. Machine-learning forecasters (gradient boosting on lagged features, LSTMs) add external drivers such as operational tempo. Whatever the method, evaluate forecasts on later data only (\emph{rolling-origin} evaluation) --- random splits of a time series leak the future into the past.

\section{Non-linear Tracking in Detail}
\subsection{The matrix Kalman filter as Bayesian updating}
With state $x_k=Fx_{k-1}+w_k$, $w_k\sim\mathcal N(0,Q)$, and measurement $z_k=Hx_k+v_k$, $v_k\sim\mathcal N(0,R)$, a Gaussian prior propagated through linear models stays Gaussian. The predict step computes the prior $\mathcal N(F\hat x,FPF^\top+Q)$; the update step is Bayes' rule for two Gaussians, whose product is Gaussian with
\[K=P^-H^\top(HP^-H^\top+R)^{-1},\quad\hat x=\hat x^-+K(z-H\hat x^-),\quad P=(I-KH)P^-.\]
The innovation $\nu=z-H\hat x^-$ has covariance $S=HP^-H^\top+R$; its consistency with $S$ is the filter's health check (Section~\ref{sec:gating}).
\subsection{The extended Kalman filter: linearise}
Radars measure range and bearing, non-linear functions of Cartesian position: $h(x,y)=\big(\sqrt{x^2+y^2},\;\operatorname{atan2}(y,x)\big)$. The EKF replaces $H$ by the Jacobian of $h$ at the current estimate.
\begin{example}{a range--bearing Jacobian}
At estimated position $(3{,}000,\,4{,}000)$ m the range is 5,000 m and the bearing $53.13^\circ$. The Jacobian is
\[H=\begin{pmatrix}x/r&y/r\\-y/r^2&x/r^2\end{pmatrix}=\begin{pmatrix}0.6&0.8\\-1.6\times10^{-4}&1.2\times10^{-4}\end{pmatrix}.\]
The first row says a 1 m move along the line of sight changes the range by 1 m; the second says a 1 m move across it changes the bearing by only $2\times10^{-4}$ rad. The EKF is accurate when the function is nearly linear over the uncertainty region; with large uncertainty or strong curvature it can diverge.
\end{example}
\subsection{The unscented Kalman filter: sample}
The UKF passes a few deterministically chosen \textbf{sigma points} through the true non-linear function and recomputes the mean and covariance from them.
\begin{example}{sigma points for $y=x^2$}
Let $x\sim\mathcal N(1,0.5^2)$. Linearising at the mean (EKF) predicts $E[y]=1$. The true mean is $E[x^2]=1+0.25=1.25$. With one dimension and $\kappa=2$, the sigma points are $1$ and $1\pm\sqrt3\times0.5=1.866,\,0.134$, with weights $2/3$, $1/6$, $1/6$. Their transformed mean is $\frac23(1)+\frac16(3.482)+\frac16(0.018)=1.25$, and their variance $1.125$, both exactly correct here. The EKF's linearisation missed the curvature that the sigma points captured.
\end{example}
\subsection{Particle filters: simulate}
When the distribution is far from Gaussian --- multiple hypotheses about where a target went behind a ridge --- a \textbf{particle filter} represents it by many weighted samples. Each cycle: move every particle through the motion model with random noise; weight each by the likelihood of the measurement; resample particles in proportion to their weights.
\begin{example}{three particles}
Particles at 48, 51 and 55 km; a measurement of 52 km with variance 4. Likelihood weights $\propto\exp(-(52-x)^2/8)$: 0.135, 0.882, 0.325, normalised to 0.10, 0.66, 0.24. Weighted estimate 51.7 km. Resampling would copy the 51-km particle most often. In practice thousands of particles are used.
\end{example}
\subsection{Manoeuvring targets}
A single motion model cannot describe a target that cruises and then turns hard. The \textbf{interacting multiple model} (IMM) filter runs several filters in parallel (constant velocity, coordinated turn, constant acceleration) and mixes them with model probabilities updated by how well each predicts the measurements.

\section{Fusing Evidence and Associating Tracks}
\subsection{Levels of fusion}
The JDL data-fusion model distinguishes: level~0, signal refinement; level~1, \textbf{object assessment} (tracking and identification --- this lecture); level~2, \textbf{situation assessment} (relationships among objects: a formation, an attack); level~3, \textbf{impact assessment} (threat and consequences); level~4, process refinement (tasking sensors to reduce uncertainty). Machine learning contributes at every level, but most mature systems are at level~1.

\subsection{Dempster--Shafer evidence combination}
Probability forces belief to be split between ``hostile'' and ``friendly'' even when evidence says nothing. Dempster--Shafer theory lets a source assign mass to the uncertain set $\Theta=\{H,F\}$ (``don't know'') and combines sources by Dempster's rule: multiply masses, assign each product to the intersection of the two sets, discard products with empty intersection (the \textbf{conflict} $K$), and renormalise by $1-K$.
\begin{example}{combining radar and ESM evidence}
Radar: $m_1(H)=0.6$, $m_1(F)=0.1$, $m_1(\Theta)=0.3$. ESM: $m_2(H)=0.5$, $m_2(F)=0.3$, $m_2(\Theta)=0.2$. Products to $H$: $0.30+0.12+0.15=0.57$; to $F$: $0.03+0.02+0.09=0.14$; to $\Theta$: $0.06$; conflict ($H$ with $F$): $0.18+0.05=0.23$. Normalised: $m(H)=0.74$, $m(F)=0.18$, $m(\Theta)=0.08$. Agreement on ``hostile'' is reinforced and the ``don't know'' mass shrinks. A large conflict $K$ is itself a warning that the sources disagree --- possibly because one is being spoofed --- and Dempster's rule behaves badly when $K$ approaches 1.
\end{example}

\subsection{Track-to-measurement association and gating}
\label{sec:gating}
Before a measurement can update a track, the system must decide which track it belongs to. \textbf{Gating} discards implausible pairs using the squared Mahalanobis distance of the innovation, $d^2=\nu^\top S^{-1}\nu$, which follows a $\chi^2$ distribution with as many degrees of freedom as the measurement has dimensions.
\begin{example}{a gate test}
Innovation $\nu=(3,-2)$ with $S=\text{diag}(4,1)$: $d^2=9/4+4/1=6.25$. The 99\% gate for two dimensions is $\chi^2_{2,0.99}=9.21$, so the measurement is inside the gate and is a candidate for this track; at the 95\% gate (5.99) it would be rejected. Gate size trades missed associations against clutter.
\end{example}
Among the candidates, \textbf{nearest neighbour} takes the closest; \textbf{global nearest neighbour} solves an assignment problem across all tracks and measurements at once (the Hungarian algorithm); \textbf{JPDA} averages over all feasible assignments weighted by their probabilities; \textbf{multiple hypothesis tracking} defers hard decisions and keeps alternative association histories alive until later data resolve them. In dense or deceptive environments --- decoys, chaff, crossing targets --- association errors, not filter errors, dominate tracking failures.

\section{Further Self-Test Questions}
\begin{question}{Part C}
(1) Why is a prediction interval always wider than the confidence interval for the mean? (2) A Poisson model has $\beta_1=-0.4$ for ``armoured vehicles per platoon''. Interpret it. (3) Under the square law, Red's kill rate is 1.5 times Blue's. What force ratio does Blue need? (4) Combine $m_1(H)=0.8,m_1(\Theta)=0.2$ with $m_2(F)=0.8,m_2(\Theta)=0.2$ by Dempster's rule. (5) Why must time-series forecasts be evaluated on later data only?\\
\textit{Answers.} (1) It adds the irreducible noise of an individual outcome ($\sigma^2$) to the uncertainty of the fitted line. (2) Each additional armoured vehicle multiplies expected casualties by $e^{-0.4}=0.67$, a 33\% reduction --- an association in the data, not necessarily a causal effect. (3) $\sqrt{1.5}=1.22$. (4) Conflict $K=0.64$; $m(H)=0.16/0.36=0.44$, $m(F)=0.44$, $m(\Theta)=0.04/0.36=0.11$: strong conflict, flag the sources. (5) Random splits let the model learn from the future, inflating accuracy.
\end{question}

\section{Implementation Notes}
\begin{center}\begin{minipage}{0.92\textwidth}\small\ttfamily
from sklearn.linear\_model import RidgeCV, LassoCV, ElasticNetCV\\
from xgboost import XGBRegressor\\[2pt]
ridge = RidgeCV(alphas=np.logspace(-3, 3, 25)).fit(Xs\_train, y\_train)\\
xgb = XGBRegressor(n\_estimators=2000, learning\_rate=0.05, max\_depth=5,\\
\hspace*{4em}subsample=0.8, colsample\_bytree=0.8)\\
xgb.fit(X\_train, y\_train, eval\_set=[(X\_val, y\_val)], early\_stopping\_rounds=50)\\
q90 = XGBRegressor(objective="reg:quantileerror", quantile\_alpha=0.9)
\end{minipage}\end{center}
Early stopping finds the number of trees automatically (about 400 here); the quantile model gives planning margins.

\section{Self-Test Questions}
\begin{question}{five questions}
\textbf{1.} A residual plot widens with fitted value. Name the problem and two remedies. \textbf{2.} Why must features be standardised before ridge or lasso? \textbf{3.} In the 2-D Kalman example, why did the velocity estimate change although only position was measured? \textbf{4.} What does a normalised innovation squared persistently near 4 (for 1-D measurements) indicate? \textbf{5.} Why is naive information-form fusion dangerous for two tracks that share a past measurement?

\textit{Answers.} 1: heteroscedasticity; log-transform the target, weighted or quantile regression. 2: the penalty would depend on arbitrary units. 3: the predicted covariance has non-zero position--velocity terms, so the gain has a velocity component. 4: the filter is over-confident ($Q$ or $R$ too small) or the motion model is wrong. 5: it double-counts correlated information and under-states uncertainty; use covariance intersection.
\end{question}

\section{Regression Across the Disciplines}
\begin{itemize}
\item \textbf{Biology/statistics:} Galton (1886) named ``regression'' after regression towards the mean in heights.
\item \textbf{Aeronautics:} the Kalman filter was developed for Apollo navigation at NASA Ames in the early 1960s.
\item \textbf{Engineering:} fatigue S--N curves regress log-cycles-to-failure on stress to set component lives.
\item \textbf{Psychology:} the Yerkes--Dodson inverted-U (1908) relating arousal to performance cannot be fitted by a straight line.
\item \textbf{Operations research:} Engel (1954) fitted Lanchester's square law to the Iwo Jima battle record.
\item \textbf{Toxicology:} dose--response (logistic/probit) regression underpins CBRN casualty estimates.
\end{itemize}

\section{Summary and Glossary}
\begin{center}\small\renewcommand{\arraystretch}{1.2}
\begin{tabular}{lccccc}\toprule Case A model & $R^2$ & RMSE & MAE & Train (s) & Features\\\midrule OLS & 0.710 & 4.22 & 3.12 & 0.06 & 18\\ Ridge & 0.738 & 3.91 & 2.88 & 0.07 & 18\\ Lasso & 0.728 & 4.02 & 2.94 & 0.14 & 12\\ Elastic net & 0.741 & 3.88 & 2.81 & 0.18 & 10\\ SVR & 0.812 & 3.24 & 2.37 & 4.8 & 18\\ \textbf{XGBoost} & \textbf{0.912} & \textbf{2.06} & \textbf{1.54} & 22.4 & 18\\\bottomrule\end{tabular}\end{center}
\begin{center}\small\renewcommand{\arraystretch}{1.2}
\begin{tabularx}{\textwidth}{>{\bfseries}p{3.3cm}X}\toprule
Aleatoric uncertainty & Irreducible randomness in the outcome.\\
Epistemic uncertainty & Uncertainty from limited data or knowledge.\\
Innovation & Difference between a measurement and its predicted value.\\
Kalman gain & Weight given to the innovation; $P/(P+R)$ in one dimension.\\
Multicollinearity & Near-linear dependence among features, inflating coefficient variance.\\
Shrinkage & Pulling coefficients toward zero to reduce variance.\\
Soft-thresholding & Lasso's rule setting small coefficients exactly to zero.\\
VIF & $1/(1-R_j^2)$: inflation of a coefficient's variance from collinearity.\\\bottomrule
\end{tabularx}\end{center}
\end{document}
